\textbf{RH1 (Absence de supériorité générale).} Sur 20 décision(s) : 11 conforme(s) à l'hypothèse, 0 contraire(s), 3 non tranchée(s), 6 différence(s) démontrée(s) sans sens prédit. HS1-sep4 : écart moyen A $-$ B de 0,003 (IC à 95 \% [$8{,}54 \times 10^{-4}$ ; 0,005], n = 110, p de Holm $8{,}67 \times 10^{-70}$), décision « équivalent » [conforme] ; HS1-sep5 : écart moyen A $-$ B de $-1{,}05 \times 10^{-4}$ (IC à 95 \% [-0,001 ; $8{,}41 \times 10^{-4}$], n = 110, p de Holm $1{,}95 \times 10^{-108}$), décision « équivalent » [conforme] ; HS1-sep1 : écart moyen A $-$ B de 0,015 (IC à 95 \% [0,011 ; 0,021], n = 110, p de Holm $5{,}20 \times 10^{-9}$), décision « A > B » [différence non orientée] ; HS1-sep2 : écart moyen A $-$ B de 0,173 (IC à 95 \% [0,154 ; 0,191], n = 110, p de Holm $1{,}26 \times 10^{-29}$), décision « A > B » [différence non orientée] ; HS1-sep3 : écart moyen A $-$ B de 0,281 (IC à 95 \% [0,264 ; 0,299], n = 110, p de Holm $3{,}67 \times 10^{-31}$), décision « A > B » [différence non orientée] ; HS2 : écart moyen A $-$ B de -0,255 (IC à 95 \% [-0,336 ; -0,173], n = 110, p de Holm $7{,}48 \times 10^{-8}$), décision « A < B » [différence non orientée] ; HS4 : écart moyen A $-$ B de 0,001 (IC à 95 \% [-0,002 ; 0,004], n = 110, p de Holm 1,0000), décision « différence non démontrée » [non tranché] ; HS5-D2 : écart moyen A $-$ B de $5{,}67 \times 10^{-6}$ (IC à 95 \% [$-3{,}48 \times 10^{-4}$ ; $3{,}61 \times 10^{-4}$], n = 110, p de Holm 1,0000), décision « différence non démontrée » [non tranché] ; HS5-D5 : écart moyen A $-$ B de $-5{,}34 \times 10^{-4}$ (IC à 95 \% [-0,001 ; $-4{,}17 \times 10^{-5}$], n = 110, p de Holm 0,4150), décision « différence non démontrée » [non tranché] ; HS5-D10 : écart moyen A $-$ B de 0,282 (IC à 95 \% [0,254 ; 0,312], n = 110, p de Holm $5{,}55 \times 10^{-32}$), décision « A > B » [différence non orientée] ; HS5-D20 : écart moyen A $-$ B de 0,794 (IC à 95 \% [0,787 ; 0,801], n = 110, p de Holm $2{,}93 \times 10^{-32}$), décision « A > B » [différence non orientée] ; HS7-p0.01-a5.0 : écart moyen A $-$ B de $3{,}41 \times 10^{-4}$ (IC à 95 \% [$-5{,}60 \times 10^{-4}$ ; 0,001], n = 110, p de Holm $3{,}58 \times 10^{-113}$), décision « équivalent » [conforme] ; HS7-p0.01-a10.0 : écart moyen A $-$ B de $5{,}70 \times 10^{-4}$ (IC à 95 \% [-0,001 ; 0,002], n = 110, p de Holm $1{,}63 \times 10^{-84}$), décision « équivalent » [conforme] ; HS7-p0.01-a20.0 : écart moyen A $-$ B de $1{,}89 \times 10^{-4}$ (IC à 95 \% [-0,001 ; 0,001], n = 110, p de Holm $7{,}00 \times 10^{-98}$), décision « équivalent » [conforme] ; HS7-p0.05-a5.0 : écart moyen A $-$ B de $-1{,}74 \times 10^{-4}$ (IC à 95 \% [-0,002 ; 0,001], n = 110, p de Holm $2{,}30 \times 10^{-87}$), décision « équivalent » [conforme] ; HS7-p0.05-a10.0 : écart moyen A $-$ B de $-3{,}43 \times 10^{-4}$ (IC à 95 \% [-0,001 ; $5{,}98 \times 10^{-4}$], n = 110, p de Holm $3{,}77 \times 10^{-107}$), décision « équivalent » [conforme] ; HS7-p0.05-a20.0 : écart moyen A $-$ B de 0,015 (IC à 95 \% [0,003 ; 0,031], n = 110, p de Holm $3{,}87 \times 10^{-5}$), décision « équivalent » [conforme] ; HS7-p0.1-a5.0 : écart moyen A $-$ B de -0,002 (IC à 95 \% [-0,003 ; $-1{,}62 \times 10^{-4}$], n = 110, p de Holm $2{,}62 \times 10^{-91}$), décision « équivalent » [conforme] ; HS7-p0.1-a10.0 : écart moyen A $-$ B de $9{,}40 \times 10^{-4}$ (IC à 95 \% [$-6{,}06 \times 10^{-4}$ ; 0,002], n = 110, p de Holm $1{,}41 \times 10^{-85}$), décision « équivalent » [conforme] ; HS7-p0.1-a20.0 : écart moyen A $-$ B de 0,013 (IC à 95 \% [0,002 ; 0,026], n = 110, p de Holm $7{,}48 \times 10^{-8}$), décision « équivalent » [conforme].

\textbf{RH2 (Déplacement de la sélection de K).} Sur 3 décision(s) : 0 conforme(s) à l'hypothèse, 0 contraire(s), 3 non tranchée(s). HS3-K4 : écart moyen A $-$ B de 0,050 (IC à 95 \% [0,036 ; 0,063], n = 110, p de Holm 1,0000), décision « équivalence non démontrée » [non tranché] ; HS3-K5 : écart moyen A $-$ B de 0,041 (IC à 95 \% [0,029 ; 0,052], n = 110, p de Holm 0,3889), décision « équivalence non démontrée » [non tranché] ; HS3-K6 : écart moyen A $-$ B de 0,093 (IC à 95 \% [0,084 ; 0,103], n = 110, p de Holm 1,0000), décision « équivalence non démontrée » [non tranché].

\textbf{RH4 (Régularisation par les priors).} Sur 5 décision(s) : 5 conforme(s) à l'hypothèse, 0 contraire(s), 0 non tranchée(s). HS6-N30 : écart moyen A $-$ B de -0,764 (IC à 95 \% [-0,836 ; -0,673], n = 110, p de Holm $1{,}55 \times 10^{-24}$), décision « A < B » [conforme] ; HS6-N60 : écart moyen A $-$ B de -0,482 (IC à 95 \% [-0,573 ; -0,382], n = 110, p de Holm $3{,}11 \times 10^{-15}$), décision « A < B » [conforme] ; HS6-N120 : écart moyen A $-$ B de -0,291 (IC à 95 \% [-0,382 ; -0,209], n = 110, p de Holm $5{,}59 \times 10^{-9}$), décision « A < B » [conforme] ; HS6-N300 : écart moyen A $-$ B de -0,145 (IC à 95 \% [-0,218 ; -0,073], n = 110, p de Holm 0,0010), décision « A < B » [conforme] ; HS6-N1000 : écart moyen A $-$ B de -0,145 (IC à 95 \% [-0,218 ; -0,082], n = 110, p de Holm 0,0002), décision « A < B » [conforme].
